\chapter{Conclusions and Future Work}
\label{ch:conclusions}

This thesis set out to evaluate Feature Attribution Explanation methods under
multiple, potentially conflicting effectiveness criteria, treating the
evaluation pipeline---rather than any new explanation method or predictive
model---as the unit of contribution. The motivating problem is well documented:
the proliferation of XAI evaluation metrics has not produced consensus, because
the metrics disagree, partially duplicate one another, and are themselves of
uncertain reliability. The response developed here is a single, reproducible
pipeline that prunes redundant metrics, discounts unreliable ones, and
aggregates the survivors into an effectiveness index with explicit, auditable
configuration. This chapter summarises the outcome against each of the four
contributions and outlines the directions that the work opens.

\promoterinline{This chapter is a scaffold mapping each contribution to a
conclusion stub. The stubs that depend on the empirical verdict are to be
completed once the full 600-image run is available; no specific quantitative
outcome is asserted here.}

%% ========================================================================
\section{Summary of Contributions}
\label{sec:concl-summary}

The four contributions map onto what the thesis set out to do as follows. Each
paragraph states the methodological outcome---what was developed, proposed, or
validated---and reserves the quantitative verdict for the marked stub, to be
filled once Chapter~\ref{ch:experiments} is complete.

\paragraph{Contribution~1 --- Non-redundancy analysis of FAE evaluation
metrics.}
The thesis develops a correlation-based redundancy analysis across the six
Quantus categories (Section~\ref{sec:exp-redundancy}), with a documented
pruning rule, and derives a reduced, non-redundant metric set $M^{*}$. The
methodological outcome is a transferable procedure for diagnosing and removing
metric redundancy before any comparison is made, directly addressing the
``many correlated metrics overstate evidence'' problem identified by
\citet{tomsett2020sanity} and \citet{kadir2024evaluating}.
% TODO: fill after 600-image run
\emph{[Conclusion stub: state, after the full run, which categories proved
redundant and which carried independent signal, and how stable $M^{*}$ was
across the two models.]} \textbf{[pending full run]}

\paragraph{Contribution~2 --- Multi-criteria aggregation framework.}
The thesis proposes and implements a framework that integrates MetaQuantus
meta-validation \citep{hedstrom2024metaquantus} with the Autoweighted
aggregation of NormEnsembleXAI \citep{hryniewska2024normensemblexai}, and
introduces a novel MetaQuantus-discounted weighting operator
(Equation~\eqref{eq:mq-discount}) in which each metric's aggregation weight is
attenuated by its measured reliability. The methodological outcome is a single
weight-derivation step that fuses redundancy-aware selection with
reliability-aware weighting---the central technical novelty of the thesis,
which stands independently of the rankings it produces.
% TODO: fill after 600-image run
\emph{[Conclusion stub: state whether the reliability discount measurably
changed the metric weights and the resulting ranking.]} \textbf{[pending full
run]}

\paragraph{Contribution~3 --- Empirical characterisation of metric ensembling.}
The thesis uses the framework to characterise where aggregation of metric
scores produces rankings that diverge from single-metric rankings and where it
does not, and to compare individual against NormEnsembleXAI-ensembled
attributions (Sections~\ref{sec:exp-ranking}--\ref{sec:exp-ensemble}). The
methodological outcome is an experimental protocol for isolating the effect of
aggregation, and of attribution ensembling, on method ordering.
% TODO: fill after 600-image run
\emph{[Conclusion stub: state where aggregation diverged from single-metric
rankings, which metric families drove the divergence, and whether ensembling
improved the effectiveness index and on which axes.]} \textbf{[pending full
run]}

\paragraph{Contribution~4 --- Empirical validation on ISIC~2017.}
The thesis validates the framework on expert-annotated dermoscopic data, using
the segmentation masks to make the Localization category measurable, and
applies Friedman, Nemenyi, and Wilcoxon testing to assess the significance of
method and ensembling differences (Section~\ref{sec:exp-stats}). The
methodological outcome is an end-to-end, statistically grounded instantiation of
the framework on a real, mask-annotated medical task.
% TODO: fill after 600-image run
\emph{[Conclusion stub: state whether method differences and ensembling gains
were statistically significant, and at what effect size.]} \textbf{[pending
full run]}

Taken together, the four contributions answer the thesis objective not with a
single ``best'' FAE method---which would be dataset- and configuration-specific
and of limited durability---but with a principled, reproducible apparatus for
reaching such a verdict and for stating the conditions under which it holds.

%% ========================================================================
\section{Future Work}
\label{sec:concl-future}

The framework was deliberately bounded in scope (Chapter~\ref{ch:discussion}),
and each boundary suggests a concrete extension. The directions below are
ordered roughly from those that consolidate the present results to those that
broaden the framework's reach.

\paragraph{Scaling the statistical analysis.}
The redundancy and reliability estimates rest on a finite sample from a single
challenge distribution, and the pilot $n$ is too small for firm claims. The
immediate next step is to re-run the redundancy and meta-validation analyses on
the full 600-image test set and, ideally, on repeated subsamples, reporting
confidence intervals on the correlation and reliability estimates and a
stability analysis of the membership of $M^{*}$ and of the metric weights. This
would move the metric structure from a property of one draw of images to a
property of the domain. A related refinement is to resolve the greedy
cross-category pruning behaviour and report how a strictly within-category
pruning policy (Decision~D4) changes $M^{*}$ and the downstream ranking.

\paragraph{Automated meta-validation.}
Meta-validation is currently a configured stage applied once to a fixed metric
set. A natural extension is to fold it into a fully automated loop that, given a
new dataset or model, re-derives the reliable, non-redundant metric subset and
the MQ-discounted weights without manual intervention. This would turn the
pipeline from a study-specific configuration into a reusable
``evaluate-the-evaluators'' service, and would let the framework adapt its
metric set to each new domain rather than assuming the ISIC~2017 structure
transfers.

\paragraph{More datasets and architectures.}
All evidence here comes from ISIC~2017 dermoscopy and from two convolutional
classifiers. Extending the study to a second mask-annotated medical dataset
would test whether the redundancy and reliability structure transfers across
domains, while adding a transformer-based classifier would test whether
gradient-based attributions and the layer-wise MPRT behave differently outside
the convolutional regime. Both extensions directly probe the external-validity
limits identified in Chapter~\ref{ch:discussion}.

\paragraph{Additional FAE methods, including the deferred perturbation
methods.}
The evaluated panel was restricted to seven methods for reasons of
computational cost and metric compatibility; LIME and KernelSHAP were deferred
rather than dismissed. Re-introducing the perturbation-based methods---and
resolving their incompatibility with the differentiable Robustness metrics,
perhaps via surrogate-gradient or finite-difference approximations---would test
whether the framework's conclusions hold across the gradient/perturbation
divide and would let the redundancy analysis speak to whether perturbation-based
faithfulness metrics are redundant with their gradient-based counterparts.

\paragraph{Aggregation operators and weighting.}
The thesis recommends the weighted arithmetic mean but reports the minimum and
geometric-mean operators as alternatives encoding worst-dimension and
catastrophic-failure-averse stances respectively. A focused study could compare
these operators under the MQ-discount weighting, particularly for
safety-critical interpretations where the minimum operator's
``acceptable-on-every-axis'' semantics are most defensible, and could explore
learned or preference-elicited weightings as a successor to the data-adaptive
Autoweighted scheme.

\paragraph{Regulatory alignment and auditability.}
Finally, the reproducible, meta-validated pipeline is well positioned to serve
as a transparency instrument under emerging regulation. Following
\citet{hedstrom2025bridging}, a promising direction is to connect the framework
explicitly to the transparency and documentation requirements of the EU AI Act,
formalising the pipeline's configuration record as an auditable evaluation
artefact and assessing what additional evidence---human-grounded validation in
particular---a regulator would require before a functionally grounded
effectiveness index could support a compliance claim. This would extend the
contribution from a research methodology toward a candidate evaluation standard.
